\documentclass[11pt]{article}
% ===========================================================================
%  Shared preamble for the AP Calculus AB exam-preparation worksheets.
%  Input right after \documentclass. Files beginning with "_" are not compiled
%  on their own. Navy and gold, to match the course's lesson pages.
% ===========================================================================
\usepackage[letterpaper,top=2.3cm,bottom=2.2cm,left=1.7cm,right=1.7cm,headheight=28pt,footskip=30pt]{geometry}
\usepackage{amsmath,amssymb}
\usepackage[T1]{fontenc}
\usepackage[utf8]{inputenc}
\usepackage{lmodern}
\usepackage{xcolor}
\usepackage[most]{tcolorbox}
\usepackage{enumitem}
\usepackage{fancyhdr}
\usepackage{array}
\usepackage{tikz}
\usepackage{pgfplots}
\pgfplotsset{compat=1.18}

\definecolor{navy}{HTML}{161E54}
\definecolor{ink}{HTML}{27357E}
\definecolor{gold}{HTML}{B98A00}
\definecolor{goldbg}{HTML}{FFFDF3}
\definecolor{inkbg}{HTML}{F6F7FD}
\definecolor{rule}{HTML}{C9CFF0}
\definecolor{okgreen}{HTML}{1B7A44}
\definecolor{okbg}{HTML}{F1F8F3}

\pagestyle{fancy}
\fancyhf{}
\fancyhead[L]{\footnotesize Name: \underline{\hspace{4cm}}}
\fancyhead[C]{\textbf{\color{navy}Integration Academy}\\[-2pt]{\scriptsize Calculus AB \textperiodcentered{} AP\textsuperscript{\textregistered} Exam Preparation}}
\fancyfoot[L]{\tiny\color{black!55}\parbox[t]{0.78\textwidth}{AP\textsuperscript{\textregistered} is a registered trademark of the College Board, which was not involved in the production of, and does not endorse, this material. \textcopyright{} Integration Academy.}}
\fancyfoot[R]{\scriptsize\thepage}
\renewcommand{\headrulewidth}{1.2pt}
\renewcommand{\headrule}{\hbox to\headwidth{\color{gold}\leaders\hrule height \headrulewidth\hfill}}
\renewcommand{\footrulewidth}{0pt}
\setlength{\parindent}{0pt}
\setlength{\parskip}{4pt}

% Title block: \wstitle{code}{title}{unit line}
\newcommand{\wstitle}[3]{\begin{center}{\small\bfseries\color{gold}WORKSHEET #1}\\[3pt]{\LARGE\bfseries\color{navy}#2}\\[3pt]{\small\color{black!60}#3}\end{center}\vspace{2pt}}
% Section heading
\newcommand{\sech}[1]{\par\vspace{10pt}{\large\bfseries\color{navy}#1}\par\nobreak\vspace{2pt}{\color{rule}\hrule height0.8pt}\vspace{7pt}}
% A level tag printed beside a group of questions
\newcommand{\level}[1]{\par\vspace{6pt}{\footnotesize\bfseries\color{gold}\MakeUppercase{#1}}\par\nobreak\vspace{2pt}}

\newtcolorbox{ideabox}{colframe=ink,colback=inkbg,coltitle=white,colbacktitle=ink,fonttitle=\bfseries\small,title={Key ideas},arc=2pt,boxrule=0.8pt,left=8pt,right=8pt}
\newtcolorbox{example}[1]{colframe=rule,colback=white,coltitle=navy,colbacktitle=inkbg,fonttitle=\bfseries\small,title={#1},arc=2pt,boxrule=0.8pt,borderline west={3pt}{0pt}{ink},left=10pt,right=8pt}
\newtcolorbox{frbox}[1]{colframe=gold,colback=goldbg,coltitle=white,colbacktitle=gold,fonttitle=\bfseries\small,title={#1},arc=2pt,boxrule=0.8pt,breakable,left=8pt,right=8pt}
\newtcolorbox{solbox}{colframe=okgreen,colback=okbg,coltitle=white,colbacktitle=okgreen,fonttitle=\bfseries\small,title={Worked solutions},arc=2pt,boxrule=0.8pt,breakable,left=8pt,right=8pt}
\newtcolorbox{scorebox}[1]{colframe=navy,colback=white,coltitle=white,colbacktitle=navy,fonttitle=\bfseries\small,title={#1},arc=2pt,boxrule=0.8pt,breakable,left=8pt,right=8pt}

\newcommand{\soln}{\par\smallskip\textbf{\color{ink}Solution.}\ }
\newcommand{\result}[1]{\par\smallskip\textbf{#1}}
% blank working space of a given height
\newcommand{\work}[1]{\par\vspace{#1}}
% a point value in the margin of a scoring guide
\newcommand{\pts}[1]{\hfill{\footnotesize\bfseries\color{gold}[#1]}}

% Coordinate grid for graph questions: \begin{apgraph}{xmin}{xmax}{ymin}{ymax} ... \end{apgraph}
\newenvironment{apgraph}[4]{\begin{center}\begin{tikzpicture}\begin{axis}[
  axis lines=middle, xmin=#1, xmax=#2, ymin=#3, ymax=#4,
  grid=both, grid style={black!12}, xtick distance=1, ytick distance=1,
  ticklabel style={font=\scriptsize}, xlabel={$x$}, ylabel={$y$},
  width=10cm, height=7.2cm, samples=80, clip=false,
  every axis plot/.append style={line width=1.1pt, color=ink}]}{\end{axis}\end{tikzpicture}\end{center}}
\newcommand{\closeddot}[2]{\addplot[only marks, mark=*, mark size=2.2pt, color=ink] coordinates {(#1,#2)};}
\newcommand{\opendot}[2]{\addplot[only marks, mark=*, mark size=2.2pt, color=ink, mark options={fill=white}] coordinates {(#1,#2)};}

% Multiple-choice options: four across, or stacked when the options are long.
\newcommand{\opts}[4]{\par\vspace{3pt}\begin{tabular}{@{}p{0.235\linewidth}p{0.235\linewidth}p{0.235\linewidth}p{0.235\linewidth}@{}}(A)\ #1 & (B)\ #2 & (C)\ #3 & (D)\ #4\end{tabular}\par\vspace{9pt}}
\newcommand{\optsv}[4]{\par\vspace{3pt}\begin{tabular}{@{}l@{\ \ }p{0.9\linewidth}@{}}(A) & #1\\(B) & #2\\(C) & #3\\(D) & #4\end{tabular}\par\vspace{9pt}}

\fancyhead[R]{\footnotesize\color{black!60}Worksheet 1.6}
\begin{document}

\wstitle{1.6}{Trapping a Limit: The Squeeze Theorem}{Unit 1: Limits and Continuity}

\begin{ideabox}
\begin{itemize}[leftmargin=*,itemsep=2pt,topsep=2pt]
  \item \textbf{Squeeze theorem.} If $g(x)\le f(x)\le h(x)$ near $c$ and $\displaystyle\lim_{x\to c}g(x)=\displaystyle\lim_{x\to c}h(x)=L$, then $\displaystyle\lim_{x\to c}f(x)=L$.
  \item Bounds usually come from $-1\le\sin(\cdot)\le1$ or $-1\le\cos(\cdot)\le1$, multiplied by a non-negative factor.
  \item $\displaystyle\lim_{\theta\to 0}\dfrac{\sin\theta}{\theta}=1$ and $\displaystyle\lim_{\theta\to 0}\dfrac{1-\cos\theta}{\theta}=0$, with $\theta$ in radians.
  \item The expression inside the sine must match the expression underneath it.
\end{itemize}
\end{ideabox}

\sech{Worked examples}
\begin{example}{Example 1 \textperiodcentered{} Matching the argument}
Evaluate $\displaystyle\lim_{x\to 0}\dfrac{\sin(7x)}{2x}$.\soln $\dfrac{\sin(7x)}{2x}=\dfrac72\cdot\dfrac{\sin(7x)}{7x}\to\dfrac72\cdot1$.\result{The limit is $\dfrac72$.}
\end{example}

\begin{example}{Example 2 \textperiodcentered{} A function known only by its bounds}
A function satisfies $2x+1\le f(x)\le x^2+2$ for all $x$. Find $\displaystyle\lim_{x\to 1}f(x)$.\soln $\displaystyle\lim_{x\to 1}(2x+1)=3$ and $\displaystyle\lim_{x\to 1}(x^2+2)=3$.\result{By the squeeze theorem, $\displaystyle\lim_{x\to 1}f(x)=3$.}
\end{example}

\begin{example}{Example 3 \textperiodcentered{} A full squeeze argument}
Show that $\displaystyle\lim_{x\to 0}x^2\sin\!\left(\dfrac3x\right)=0$.\soln For $x\neq0$, $-1\le\sin\!\left(\tfrac3x\right)\le1$. Since $x^2\ge0$, $-x^2\le x^2\sin\!\left(\tfrac3x\right)\le x^2$. Also $\displaystyle\lim_{x\to 0}(-x^2)=0$ and $\displaystyle\lim_{x\to 0}x^2=0$.\result{By the squeeze theorem the limit is $0$.}
\end{example}

\newpage
\sech{Practice}

\level{Core}
\begin{enumerate}[leftmargin=*,itemsep=2pt]
  \item Evaluate $\displaystyle\lim_{x\to 0}\dfrac{\sin(4x)}{x}$.\work{2.3cm}
  \item Evaluate $\displaystyle\lim_{x\to 0}\dfrac{\sin x}{5x}$.\work{2.3cm}
  \item Evaluate $\displaystyle\lim_{x\to 0}\dfrac{\tan(3x)}{x}$.\work{2.3cm}
  \item Evaluate $\displaystyle\lim_{x\to 0}\dfrac{x}{\sin(2x)}$.\work{2.3cm}
\end{enumerate}

\level{Exam level}
\begin{enumerate}[leftmargin=*,itemsep=2pt,start=5]
  \item Evaluate $\displaystyle\lim_{x\to 0}\dfrac{\sin(6x)}{\sin(2x)}$.\work{2.8cm}
  \item Evaluate $\displaystyle\lim_{x\to 0}\dfrac{1-\cos(2x)}{x^2}$.\work{2.8cm}
  \item A function satisfies $5-x^2\le g(x)\le5+x^4$ for all $x$. Find $\displaystyle\lim_{x\to 0}g(x)$, with justification.\work{2.8cm}
  \item Evaluate $\displaystyle\lim_{x\to \infty}\dfrac{3+\sin x}{x}$, with justification.\work{2.8cm}
\end{enumerate}

\begin{frbox}{Question 9 \textperiodcentered{} Free-response practice}
Let $f(x)=x^3\cos\!\left(\dfrac1x\right)$ for $x\neq0$.
\begin{enumerate}[label=(\alph*),leftmargin=*,itemsep=2pt]
  \item Explain why $\displaystyle\lim_{x\to 0}f(x)$ cannot be found with the product law.\work{2.6cm}
  \item Use the squeeze theorem to find $\displaystyle\lim_{x\to 0}f(x)$. (Take care: $x^3$ can be negative.)\work{2.6cm}
  \item Does the same method show that $\displaystyle\lim_{x\to 0}(x+1)\cos\!\left(\dfrac1x\right)$ exists? Explain.\work{2.6cm}
\end{enumerate}
\end{frbox}

\newpage
\begin{solbox}
\begin{enumerate}[leftmargin=*,itemsep=4pt]
  \item $4$.
  \item $\dfrac15$.
  \item $\dfrac{\sin(3x)}{x}\cdot\dfrac{1}{\cos(3x)}\to3$.
  \item $\dfrac12\cdot\dfrac{2x}{\sin(2x)}\to\dfrac12$.
  \item $\dfrac{\sin 6x}{6x}\cdot\dfrac{2x}{\sin 2x}\cdot\dfrac{6}{2}\to3$.
  \item $\dfrac{1-\cos2x}{x^2}\cdot\dfrac{1+\cos2x}{1+\cos2x}=\dfrac{\sin^2(2x)}{x^2}\cdot\dfrac{1}{1+\cos2x}\to4\cdot\dfrac12=2$.
  \item $\displaystyle\lim_{x\to 0}(5-x^2)=5$ and $\displaystyle\lim_{x\to 0}(5+x^4)=5$, so by the squeeze theorem the limit is $5$.
  \item For $x>0$, $\dfrac2x\le\dfrac{3+\sin x}{x}\le\dfrac4x$. Both bounds tend to 0 as $x\to\infty$, so by the squeeze theorem the limit is $0$.
  \item \begin{enumerate}[label=(\alph*),leftmargin=*,itemsep=3pt]
    \item $\displaystyle\lim_{x\to 0}\cos\!\left(\tfrac1x\right)$ does not exist (the values oscillate), and the product law needs both limits to exist.
    \item $\left|\cos\tfrac1x\right|\le1$, so $-|x|^3\le x^3\cos\tfrac1x\le|x|^3$. Both bounds tend to 0 as $x\to0$, so by the squeeze theorem the limit is $0$.
    \item No. The bounds would be $-|x+1|$ and $|x+1|$, which tend to $-1$ and $1$: not equal. In fact the values oscillate between about $-1$ and $1$, and the limit does not exist.
  \end{enumerate}
\end{enumerate}
\end{solbox}

\end{document}
